%% kapa: exam=1 round=27 year=2016 period=1 form=A subject=economics subject_ko=경제학원론 seq=6 number=46
%% source: Q-Net 감정평가사 기출문제 (cst00309) artlSeq=5169402 file=extracted/1008/5169402/제27회 감정평가사 1차 문제/제27회 감정평가사 1교시문제(A형).pdf
%% conversion: kordoc 4.24.51 -> tools/gen1.py (auto); review=proofread-sonnet; flags=none
\begin{question}{46}
X재의 생산과정에서 양(+)의 외부효과가 발생할 때 균형산출량 수준에서 옳은 것은? (단, X재 시장은 완전경쟁시장이고, X재에 대한 수요의 법칙과 공급의 법칙이 성립하며, 정부의 개입은 없다고 가정한다. P는 X재의 가격, PMC는 X재의 사적 한계비용, SMC는 X재의 사회적 한계비용이다.)
\begin{choices}
  \item P = SMC = PMC
  \item P = PMC $>$ SMC
  \item P = PMC $<$ SMC
  \item P = SMC $<$ PMC
  \item PMC $<$ SMC $<$ P
\end{choices}
\end{question}
\answer{2}
